\subsection*{Vanishing of the compact quadratic cup and its actual obstruction transfer}
\begin{statement}{E.113}{Proved}{The complete compact cup calculation in the odd and rational quartic targets}
Let $E=E_+:y^2=x^3+x$.  For every prime $\ell$, including $2$,
the Weil-pairing cup map is zero:
\[
 H^1_{\rm cts}(\mathbb Q,T_\ell E)^{\otimes2}
 \xrightarrow{\ \cup\ }
 H^2_{\rm cts}(\mathbb Q,T_\ell E\otimes T_\ell E)
 \xrightarrow{\ e_*\ }
 H^2_{\rm cts}(\mathbb Q,\mathbb Z_\ell(1)),
 \qquad e_*\cup=0.
 \tag{E113.1}
\]
All cohomology here is unrestricted global continuous cohomology.
The finite-level, real-place, and inverse-limit maps proving
(E113.1) are given below.

Now take the rings $R$ and the actual quartic origin $S_R$ of
E.105: an odd $\mathbb Z_\ell$, their product, or any
$\mathbb Q_\ell$.  Put $V=T_RE$, $C=V^3$, and retain the
actual E.103 coefficient splitting
\[
 A=W\oplus j(C),\qquad j=j_*/4,
 \qquad B=\Lambda^2 A.
 \tag{E113.2}
\]
The elliptic maps and the choice of the single generator
$[h]\in H^1(\mathbb Q,V)=R$ are precisely those of E.103.
The full tensor cup on these global elliptic classes is zero.
Consequently the compact quadratic obstruction is zero for
every compact coordinate $(r_1[h],r_2[h],r_3[h])$.
This includes all six determinant coordinates and all three
$\operatorname{Sym}^2V$ coordinates of (E105.8).

Here are literal cochain and coefficient maps for that conclusion.
Choose one continuous cocycle $h$ representing $[h]$, and choose
one continuous $R(1)$-valued cochain $k$ satisfying
\[
 dk(g,t)=\langle h(g),g h(t)\rangle,
 \qquad dk(g,t)=gk(t)-k(gt)+k(g).
 \tag{E113.3}
\]
The construction in the proof supplies such a $k$ from compatible
finite primitives, and applies in every stated ring.
For $\mathbf r=(r_1,r_2,r_3)$ define
$j_{\mathbf r}(v)=j(r_1v,r_2v,r_3v)$ and
\[
 s_{\mathbf r}:R(1)\longrightarrow B,
 \qquad s_{\mathbf r}=(\Lambda^2 j_{\mathbf r})e^{-1},
 \qquad \nu_{\mathbf r}=\tfrac12s_{\mathbf r}k.
 \tag{E113.4}
\]
In a symplectic basis $(e,f)$, write
$e_i=j(0,\ldots,e,\ldots,0)$ and similarly $f_i$.
The coefficient column is exactly
\[
 s_{\mathbf r}(1)=
 \sum_i r_i^2 e_i\wedge f_i+
 \sum_{i<j}r_ir_j(e_i\wedge f_j-f_i\wedge e_j).
 \tag{E113.5}
\]
Its six determinant projection coordinates, in the order
$(11,22,33,12,13,23)$, are
$(r_1^2,r_2^2,r_3^2,2r_1r_2,2r_1r_3,2r_2r_3)$.
Its symmetric-square projection is zero.  In particular
\[
 d\nu_{\mathbf r}(g,t)
       =\tfrac12j_{\mathbf r}h(g)\wedge
                         g j_{\mathbf r}h(t).
 \tag{E113.6}
\]

The transfer to the actual remaining quartic obstruction also has
a specified inverse.  In the BCH coordinates relative to $S_R$
write the actual action as
\[
 g(x,z)=(\rho_A(g)x,\rho_B(g)z+D_gx),
 \quad c_g=D_g\rho_A(g)^{-1},
 \quad c_{gt}=c_g+g c_t.
 \tag{E113.7}
\]
The $D_g$ in this formula is the actual arithmetic operator.
For any aligned abelian cocycle $a=b+e_{\mathbf r}$, with
$b\in Z^1(\mathbb Q,W)$ and $e_{\mathbf r}=j_{\mathbf r}h$,
put
\[
 \begin{split}
 \Omega_b(g,t)&=D_gb(t)+\tfrac12b(g)\wedge g b(t),\\
 P_{b,\mathbf r}(g)&=\nu_{\mathbf r}(g)
                         +\tfrac12b(g)\wedge e_{\mathbf r}(g),\\
 \Omega^{\rm rem}_{b,\mathbf r}(g,t)
       &=\Omega_b(g,t)+D_ge_{\mathbf r}(t)
                                +b(g)\wedge g e_{\mathbf r}(t).
 \end{split}
 \tag{E113.8}
\]
Then the exact defect and its change of coordinates are
\[
 \begin{gathered}
 \Omega_a(g,t)=D_ga(t)+\tfrac12a(g)\wedge g a(t)
       =\Omega^{\rm rem}_{b,\mathbf r}(g,t)+dP_{b,\mathbf r}(g,t),\\
 (a,z)\longmapsto(b,\mathbf r,\widetilde z=z+P_{b,\mathbf r}),
 \qquad
 d\widetilde z=-\Omega^{\rm rem}_{b,\mathbf r},
 \qquad z=\widetilde z-P_{b,\mathbf r}.
 \end{gathered}
 \tag{E113.9}
\]
These maps give exactly the class-two lifts of every E.103
abelian input, with the full central-conjugacy relation.
The toric term is the one evaluated in E.108; the two additional
terms in (E113.8) retain respectively the actual handle action
and the mixed cup.  At $\ell=5$ the former transports to E.109's
nonzero compact outer class by (E113.19) below.
The computation therefore removes the entire compact quadratic
term while retaining those remaining arithmetic inputs in E.23.

There is also a complete vanishing statement for nine actual
coefficient projections of the \emph{total} obstruction.
Use E.108's direct-summand projection
$P_{\rm T}:B\to B_{\rm T}=R(1)^6\oplus(R(1)\chi_{-1})^3$
and its section $S_{\rm T}$, displayed in (E113.20).
For every actual abelian input $a$, one has
$P_{{\rm T}*}[\Omega_a]=0$.
Formula (E113.22) supplies a continuous primitive
$q_a$ with $dq_a=P_{\rm T}\Omega_a$.
Consequently $z\mapsto z+S_{\rm T}q_a$ changes the full
section equation to
$d\widetilde z=-(1-S_{\rm T}P_{\rm T})\Omega_a$.
Its $B_{\rm T}$ coordinate is an arbitrary one-cocycle.
This retains the full central fibre, including its nine
Tate-line factors; it gives the exact rank-$127$ target
for the remaining obstruction.
\end{statement}
\begin{proof}
\emph{The actual bad-prime calculation.}
The displayed integral equation has
$c_4=-48,c_6=0,\Delta=-64$.  It is minimal at $2$, since a
decrease of the discriminant valuation of an integral
Weierstrass equation is a positive multiple of $12$.
Its reduction has exactly one singular point, $(1,0)$,
and two nonsingular rational points, $O$ and $(0,0)$.
An integral point with odd $x$ would satisfy
\[
 v_2(x^3+x)=v_2\bigl(x(x^2+1)\bigr)=1,
 \tag{E113.10}
\]
which is incompatible with $y^2=x^3+x$.
Thus every integral point reduces to a nonsingular point;
every nonintegral point reduces to $O$ in the projective model.
The reduction homomorphism consequently has image of order
at most two.  The kernel is the formal group at $O$, with
$t=-x/y\in2\mathbb Z_2$.  In the chart $w=-1/y$ its defining
equation is $w=t^3+tw^2$, so the unique solution has coefficients
in $\mathbb Z_2[[t]]$.  The group law, expanded in this parameter,
is integral and has linear term $t_1+t_2$; hence
$[m](t)=mt+O(t^2)$.  For odd $m$ and nonzero
$t\in2\mathbb Z_2$ the linear term has strictly smaller valuation
than the remaining terms.  It cannot be zero.
This proves, directly on the given curve,
\[
 E(\mathbb Q_2)[\ell^n]=0
                   \quad(\ell\text{ odd},\ n\ge1).
 \tag{E113.11}
\]
The formal-group construction is also the literal one in
\cite[Chapter I, Example 2.5(c)]{E_MilneCFT}.
As a separate check, $x=X+1$ changes the equation to
$y^2=X^3+3X^2+4X+2$.  The type-II branch in
\cite[\S3.2, algorithm step 23]{E113_Cremona} applies because
$v_2(a_6)=1$.  It gives conductor exponent six and local
index one, agreeing with the direct argument and the retained
PARI calculation.

The Weil pairing identifies the Cartier dual of $E[\ell^n]$
with $E[\ell^n]$ itself.  Finite local duality
\cite[Chapter I, Theorem 2.1 and Corollary 2.3]{E_MilneADT}
therefore makes its second local cohomology the dual of its
zeroth.  Both are zero by (E113.11).  The prime-to-residue-
characteristic Euler characteristic is one; its proof is
the tame Hochschild--Serre calculation in
\cite[Chapter I, Lemma 2.9]{E_MilneADT}.
It follows that
\[
 H^1(\mathbb Q_2,E[\ell^n])=0
                   \quad(\ell\text{ odd},\ n\ge1).
 \tag{E113.12}
\]
This uses the exact finite coefficient module, so its application
to the reductions of global integral cocycles introduces no
unproved local comparison.

\emph{Good primes and the real place.}
Fix a prime $\ell$ and any two continuous $T_\ell E$-cocycles.
At $q\ne2,\ell$ the Tate module is unramified.
The restriction of either cocycle to inertia is a homomorphism.
It kills wild inertia and all prime-to-$\ell$ tame factors,
and is determined by a vector $v\in T_\ell E$ on tame
$\mathbb Z_\ell(1)$.  Arithmetic Frobenius gives
$(F_q-q)v=0$.
The elliptic Frobenius polynomial is
$X^2-a_qX+q$, and
\[
 \det(F_q-q)=q(q+1-a_q)=q\,\#E(\mathbb F_q)\ne0.
 \tag{E113.13}
\]
Thus $v=0$ in the torsion-free lattice.  The cocycles themselves,
and therefore all their finite reductions, are unramified there.
Their cup inflates from the finite-field Galois group, which
has $\ell$-cohomological dimension one.  Its second cup class
is zero.  The characteristic polynomial in (E113.13) is the
usual good-reduction Tate-module polynomial, used with the
same arithmetic-Frobenius convention in E.103.

At the real place, for odd $\ell$ the finite cohomology in
positive degree vanishes by averaging over the group of order
two.  For $\ell=2$ retain the entire compatible coefficient
tower.  Complex conjugation acts as $-1$ on $\mathbb Z_2(1)$,
and the two-periodic cyclic resolution gives
\[
 H^2(C_2,\mathbb Z_2(1))
   =\mathbb Z_2(1)^{C_2}/(1+c)\mathbb Z_2(1)=0.
 \tag{E113.14}
\]
The same conclusion for its finite reductions is visible
without passing through any completion theorem:
$H^2(\mathbb R,\mu_{2^n})$ is represented by $\{1,-1\}$,
and the map induced by squaring
$\mu_{2^{n+1}}\to\mu_{2^n}$ sends both representatives to one.
Hence every compatible family has zero component at every
finite level.  The resolution and its exactness are given in
\cite[Chapter II, Proposition 3.4 and proof]{E_MilneCFT}.

\emph{Reciprocity and the actual inverse limit.}
Let $u,v$ be the two integral cocycles, and let $w_n$ be the
Weil-pairing cup of their reductions in
$H^2(\mathbb Q,\mu_{\ell^n})$.
The Kummer sequence identifies this group with
$\operatorname{Br}(\mathbb Q)[\ell^n]$.
At every place except $\ell$ it is zero: use (E113.12) at two
when $\ell$ is odd, (E113.13) at good primes, and the preceding
real-place calculation.  The exact Brauer sequence and its
sum-of-local-invariants map
\cite[Chapter VII, Theorems 7.1, 8.1; Chapter VIII, Theorem 4.2]{E_MilneCFT}
force its remaining $\ell$-invariant to be zero, then force
$w_n=0$ globally.

The transition on
$H^1(\mathbb Q,\mu_{\ell^n})=\mathbb Q^*/\mathbb Q^{*\ell^n}$
is the surjective quotient map represented by the same rational
number.  Thus its first derived inverse limit is zero.  The
continuous cochain comparison gives an injection
\[
 H^2_{\rm cts}(\mathbb Q,\mathbb Z_\ell(1))
       \lhook\joinrel\longrightarrow
       \varprojlim_n H^2(\mathbb Q,\mu_{\ell^n}).
 \tag{E113.15}
\]
For completeness this injectivity also supplies a literal
primitive procedure.  Given a continuous two-cocycle $w$ whose
reductions are boundaries, choose $k_n$ with $dk_n=w_n$.
Suppose $k_n$ has been fixed and choose any primitive
$k'_{n+1}$.  Its reduced difference from $k_n$ is a one-cocycle.
Surjectivity of the displayed $H^1$ transition lifts its
cohomology class to a one-cocycle $t_{n+1}$.  Write, literally,
$\pi k'_{n+1}-k_n-\pi t_{n+1}=dm_n$.
For any lift $\widetilde m_n$ of that zero-cochain, set
$k_{n+1}=k'_{n+1}-t_{n+1}-d\widetilde m_n$.
Its reduction is exactly $k_n$ and its differential remains
$w_{n+1}$.
The compatible cochains define a continuous $\mathbb Z_\ell(1)$
cochain, since each finite coordinate is continuous.
This proves (E113.1) and (E113.3), without inferring an integral
vanishing solely from local or finite vanishing.

For a rational coefficient, continuous cochains on the compact
domain $G_{\mathbb Q}^m$ have bounded denominators.  Consequently
$C^m(G_{\mathbb Q},T_\ell E)\otimes\mathbb Q_\ell
=C^m(G_{\mathbb Q},V_\ell E)$, and the same statement holds
for every tensor and central coefficient used here.
Scalar extension is exact, so the integral vanishing proves
the rational one, including $\ell=2$.
For the product of odd coefficient rings, continuous maps into
the product are precisely the tuples of continuous component
maps.  The constructions and inverses apply componentwise.

\emph{All compact coefficient maps.}
By E.103 every global compact class is
$(r_1[h],r_2[h],r_3[h])$ for a single chosen generator $[h]$.
The symmetry projector on $V\otimes V$ kills $[h\cup h]$:
the degree-one cup symmetry sends this class to its negative,
and $2$ is invertible in the rings now under consideration.
The alternating projection is zero by (E113.1).
Thus the full tensor cup of any two such elliptic classes is
zero.  At the cochain level choose their common representative
$h$ and use (E113.3).

For each pair $i<j$, the actual map
$V_i\otimes V_j\to\Lambda^2C$ is
$v\otimes w\mapsto v_i\wedge w_j$.
The determinant row takes $e_i\wedge f_j$ to one and
$f_i\wedge e_j$ to minus one.  Its inverse column is
$(e_i\wedge f_j-f_i\wedge e_j)/2$.
The complementary symmetric-square columns are
$e_i\wedge e_j$,
$(e_i\wedge f_j+f_i\wedge e_j)/2$,
$f_i\wedge f_j$.
The diagonal columns are $e_i\wedge f_i$.
These give an invertible $15$-dimensional coefficient matrix
over $\mathbb Z[1/2]$, with the actual inverse retained in the
certificate.  They prove the asserted six-plus-nine decomposition
and give (E113.5) by expanding
$j_{\mathbf r}(e)\wedge j_{\mathbf r}(f)$.
The image of a common cocycle has zero symmetric-square
coordinate already as a two-cochain.  Its six determinant
coefficients are the displayed multiples of
$\langle h(g),g h(t)\rangle$.  This proves (E113.6).

\emph{The actual gauge, mixed cup, and filling maps.}
Retain the E.105 BCH product
$(x,z)(y,w)=(x+y,z+w+x\wedge y/2)$.
Block multiplication in (E113.7) proves its cocycle identity;
in particular no graded replacement of $D_g$ has been made.
For a section with geometric coordinates $(a(g),z(g))$ the
literal section equation is $dz+\Omega_a=0$.
Conjugation by a fixed geometric element $(m,w)$ is
\[
 \begin{split}
 a'(g)&=a(g)-(g-1)m,\\
 z'(g)&=z(g)+(1-g)w+\Gamma_m(g),\\
 \Gamma_m(g)&=-D_gm+\tfrac12(m+gm)\wedge a(g)
                              -\tfrac12m\wedge gm,\\
 \Omega_{a'}&=\Omega_a-d\Gamma_m.
 \end{split}
 \tag{E113.16}
\]
Thus any given compact representative can be aligned to
$e_{\mathbf r}=j_{\mathbf r}h$ using its actual abelian
coboundary and (E113.16).  The aligning compact vector is
unique because $C^{G_{\mathbb Q}}=0$; the full abelian
stabilizer vanishes by E.105.

In the coefficient submodule $W\wedge j(C)$ put
$p(g)=b(g)\wedge e_{\mathbf r}(g)$.
Expanding the two cocycle identities gives the exact equality
\[
 dp(g,t)=-b(g)\wedge g e_{\mathbf r}(t)
                         +e_{\mathbf r}(g)\wedge g b(t).
 \tag{E113.17}
\]
The mixed quadratic term is therefore
$b(g)\wedge g e_{\mathbf r}(t)+dp(g,t)/2$.
Together with (E113.6) this proves (E113.8)--(E113.9), including
both signs and both inverse coordinate maps.
The central difference of two solutions is an ordinary
one-cocycle in $B$; changing a geometric central element adds
$(1-g)w$.  Consequently the maps retain precisely the original
$H^1(\mathbb Q,B)$-torsor and all its geometric-conjugacy fibres.
They are continuous in $b,\mathbf r,z$ and in the chosen
cochain topology: only continuous linear maps, scalar products,
and the fixed cochains $h,k$ occur.  Changing the chosen
primitive $k$ changes them by its explicit central one-cocycle.

Finally let $p:A\to H=T_RJ_C$ be filling the twelve punctures,
let $Z=\Lambda^2H/Rq$ be the compact central quotient, and let
$p_B:\Lambda^2A\to Z$ be its actual exterior map.
With the chosen BCH coordinates the actual filling map is
\[
 (x,z)\longmapsto(px,p_Bz+Jx),
                   \qquad J\in\operatorname{Hom}_R(A,Z).
 \tag{E113.18}
\]
Commuting it with the actual actions gives
\[
 \begin{split}
 D^C_gp&=p_BD_g+J\rho_A(g)-\rho_Z(g)J,\\
 \Omega^C_{pa}&=p_B\Omega_a-d(Ja),\\
 c^C_g&=p_Bc_gj+Jj-g(Jj).
 \end{split}
 \tag{E113.19}
\]
Here $j:H\to A$ is the same section as in (E113.2).
Each equality follows by direct substitution in the group or
cochain law, and the middle one also verifies the actual
central-coordinate inverse on a section.
It exhibits the transfer of the surviving handle term in
(E113.8) to the compact arithmetic action.  At five, applying
E.109's coefficient left inverse $2\Pi$ to the last equality
recovers its specified nonzero CM class.  The established cup
vanishing and the actual linear arithmetic operator are thus
related by the displayed obstruction maps, with the operator
retained in every remaining-input formula.

\emph{All nine Tate-line projections of the actual action.}
The coefficient maps of E.108 give
\[
 \begin{gathered}
 P_{\rm T}:B\longrightarrow
       B_{\rm T}=R(1)^6\oplus(R(1)\chi_{-1})^3,
 \qquad S_{\rm T}:B_{\rm T}\longrightarrow B,
 \quad P_{\rm T}S_{\rm T}=1,\\
 B_{\rm rem}=\ker P_{\rm T}
   =\Lambda^2W\oplus(W\otimes C)\oplus U_{\rm CM}^3,
 \qquad \operatorname{rank}B_{\rm rem}=127.
 \end{gathered}
 \tag{E113.20}
\]
Let $\lambda$ be one row of $P_{\rm T}$, with target
$L_\lambda=R(1)$ or $R(1)\chi_{-1}$.
For the first target use $B_\lambda(x,y)=\langle x,y\rangle$.
For the second use the actual CM automorphism
$I=[i]$, $I^2=-1$, and
$B_\lambda(x,y)=\langle Ix,y\rangle$.
The identity $gIg^{-1}=\chi_{-1}(g)I$ proves its exact
equivariance.  Its symmetry follows from
$\langle Ix,y\rangle=-\langle x,Iy\rangle$.
Both perfect pairings give actual isomorphisms
$V\longrightarrow\operatorname{Hom}(V,L_\lambda)$.
The rows and sections in (E113.20) are those already
constructed over $\mathbb Z[1/2]$ in (E108.24).

For each of the three compact inclusions $j_i:V\to A$,
$\lambda c_gj_i$ is therefore a continuous cocycle in a
module explicitly isomorphic to $V$.  The complete E.103
$H^1$ calculation supplies unique scalars $s_i\in R$ and
linear maps $m_i:V\to L_\lambda$ satisfying
\[
 \begin{gathered}
 \lambda c_gj_i=s_iB_\lambda(h(g),\,\cdot\,)
                                      +(gm_i-m_i),\\
 k_\lambda=k\quad(L_\lambda=R(1)),\qquad
 k_\lambda(g)=-\tfrac12B_\lambda(h(g),h(g))
                   \quad(L_\lambda=R(1)\chi_{-1}),\\
 dk_\lambda(g,t)=B_\lambda(h(g),g h(t)).
 \end{gathered}
 \tag{E113.21}
\]
The uniqueness of $m_i$ uses the vanishing of $H^0(\mathbb Q,V)$.
For the symmetric pairing, expansion of the cocycle identity
gives the last equality directly.  Thus the arithmetic linear
term, with its actual $s_i,m_i$, has primitive
$\sum_i r_is_i k_\lambda+\sum_i r_i m_i(h)$.
This follows also by evaluating the Hom coboundary:
$d(m_i h)(g,t)=(gm_i-m_i)(g h(t))$.

For the toric input $b$, restriction of its actual defect to
$G_F$ lies in $\Lambda^2W$, by (E105.23).
It is killed by $\lambda$.  Push out the abelian-kernel
pullback over $b s_R$ along $\lambda$.
Its $G_F$-section has central coordinate zero.  Apply the
four-coset formula from the proof of (E108.14) to this actual
pushout.  The negative of the resulting section's central
coordinate is a continuous cochain $u_{\lambda,b}$ satisfying
$du_{\lambda,b}=\lambda\Omega_b$.
This construction uses division by four in precisely the
allowed rings and keeps the same abelian section.

The mixed cup lies in $W\wedge C$ and hence has zero image
under $\lambda$.  Applying (E113.9), we obtain the literal
primitive of the full projected defect:
\[
 q_{\lambda,a}
   =u_{\lambda,b}
       +\sum_i r_is_i k_\lambda
       +\sum_i r_i m_i(h)
       +\lambda P_{b,\mathbf r},
 \qquad dq_{\lambda,a}=\lambda\Omega_a.
 \tag{E113.22}
\]
For an initially unaligned representative, use (E113.16)
and subtract $\lambda\Gamma_m$ from the primitive when
passing from $a$ to $a'$.  Therefore every prescribed
representative has the asserted primitive with its exact
change-of-coordinate map.
Taking all nine rows proves (E113.20) and the final central
shift in the statement.  No value for the remaining
$U_{\rm CM}$ or mixed arithmetic operator is assigned by
this argument.
\end{proof}
